\section{Preguntas de Investigación}
\setlength{\columnseprule}{0pt}
\onecolumn
\subsection{Pregunta 1}
¿Cómo incorporar las restricciones cinemáticas no holonómicas del robot diferencial dentro de un esquema de consenso descentralizado, de forma que el protocolo de formación las respete sin recurrir a una aproximación holonómica simplificada?

\medskip
\subsection{Pregunta 2}
Dado un esquema de consenso que respeta la no holonomía del robot diferencial, ¿en qué medida una estrategia de evasión de obstáculos basada en una representación anisotrópica del entorno reduce el error de formación y el tiempo de convergencia frente a los campos de potencial isotrópicos convencionales?
